\documentclass{beamer}

\usetheme{Madrid}

\usepackage[utf8]{inputenc}
\usepackage[brazil]{babel}
\usepackage{amsmath, amssymb}
\usepackage{xcolor}

% Cor personalizada
\definecolor{mygreen}{RGB}{0,120,60}

% Aplicando cores e contrastes corretamente
\setbeamercolor{structure}{fg=mygreen}
\setbeamercolor{palette primary}{bg=mygreen, fg=white}
\setbeamercolor{palette secondary}{bg=mygreen, fg=white}
\setbeamercolor{palette tertiary}{bg=mygreen, fg=white}
\setbeamercolor{frametitle}{bg=mygreen, fg=white}
\setbeamercolor{title}{bg=mygreen, fg=white} 
\setbeamercolor{subtitle}{fg=white} % Texto em branco sobre o fundo verde da capa
\setbeamercolor{title in head/foot}{bg=mygreen, fg=white}

% Remove a barra inferior pesada do tema Madrid e deixa apenas o número do slide simples em preto
\setbeamertemplate{footline}{%
  \hfill%
  \usebeamercolor[fg]{black}%
  \small\insertframenumber%
  \hspace*{3mm}\vskip3mm%
}
\setbeamertemplate{navigation symbols}{} % Remove os ícones de navegação padrão

% ==========================================
% ORGANIZAÇÃO DOS METADADOS
% ==========================================
\title{Fatoração QR}
\subtitle{Projeto Integrador na Extensão -- Ciência da Computação} 
\author{
João Daniel \\
Kelvin Krauss \\
Sofia Effting \\
Vinicius Vincent
}
\institute{Instituto Federal Catarinense - Campus Blumenau}
\date{29 de maio de 2026}

\begin{document}

% Slide 1 - Capa
\begin{frame}
\titlepage
\end{frame}

% Slide 2
\begin{frame}{Introdução}
\begin{itemize}
    \item A fatoração QR é uma decomposição matricial.
    \item Muito utilizada em Álgebra Linear Numérica.
    \item Aplicações em computação científica e machine learning.
    \item Baseada em matrizes ortogonais.
\end{itemize}
\end{frame}

% Slide 3
\begin{frame}{Objetivos}
\begin{itemize}
    \item Compreender a fatoração QR.
    \item Entender sua construção.
    \item Relacionar QR com Gram-Schmidt.
    \item Explorar aplicações práticas.
\end{itemize}
\end{frame}

% Slide 4
\begin{frame}{Definição}
Uma matriz $A$ pode ser escrita como:

\[
A = QR
\]

Onde:

\begin{itemize}
    \item $Q$ é uma matriz ortogonal.
    \item $R$ é uma matriz triangular superior.
\end{itemize}
\end{frame}

% Slide 5
\begin{frame}{Matriz Ortogonal}
Uma matriz $Q$ é ortogonal quando:

\[
Q^TQ = I
\]

Características:
\begin{itemize}
    \item Colunas ortogonais.
    \item Norma igual a 1.
    \item Fácil inversão:
    \[
    Q^{-1} = Q^T
    \]
\end{itemize}
\end{frame}

% Slide 6
\begin{frame}{Matriz Triangular Superior}
A matriz $R$ possui formato:

\[
R =
\begin{bmatrix}
* & * & * \\
0 & * & * \\
0 & 0 & *
\end{bmatrix}
\]

Características:
\begin{itemize}
    \item Elementos abaixo da diagonal são zero.
    \item Facilita resolução de sistemas.
\end{itemize}
\end{frame}

% Slide 7
\begin{frame}{Processo de Gram-Schmidt}
Método utilizado para:
\begin{itemize}
    \item Transformar vetores independentes.
    \item Gerar base ortogonal.
    \item Construir a matriz $Q$.
\end{itemize}
\end{frame}

% Slide 8
\begin{frame}{Ortogonalização}
Dados vetores:
\[
v_1, v_2, \dots, v_n
\]

Geramos:
\[
u_1, u_2, \dots, u_n
\]

com:
\[
u_i \perp u_j
\]
\end{frame}

% Slide 9
\begin{frame}{Primeiro Vetor}
O primeiro vetor é:

\[
u_1 = v_1
\]

Normalização:

\[
q_1 = \frac{u_1}{||u_1||}
\]
\end{frame}

% Slide 10
\begin{frame}{Segundo Vetor}
Removemos projeções:

\[
u_2 = v_2 - \text{proj}_{u_1}(v_2)
\]

Onde:

\[
\text{proj}_{u_1}(v_2) =
\frac{v_2 \cdot u_1}{u_1 \cdot u_1}u_1
\]
\end{frame}

% Slide 11
\begin{frame}{Construção da Matriz Q}
A matriz $Q$ é formada pelos vetores ortonormais:

\[
Q =
\begin{bmatrix}
| & | & | \\
q_1 & q_2 & q_3 \\
| & | & |
\end{bmatrix}
\]
\end{frame}

% Slide 12
\begin{frame}{Construção da Matriz R}
A matriz $R$ é calculada por:

\[
R = Q^TA
\]

Resultado:
\begin{itemize}
    \item Matriz triangular superior.
    \item Coeficientes das projeções.
\end{itemize}
\end{frame}

% Slide 13
\begin{frame}{Exemplo Simples}
Dada a matriz $A$:

\[
A =
\begin{bmatrix}
1 & 1 \\
1 & 0
\end{bmatrix}
\]

A fatoração $A = QR$ resulta em:

\[
Q = 
\begin{bmatrix} 
\frac{\sqrt{2}}{2} & \frac{\sqrt{2}}{2} \\[6pt] 
\frac{\sqrt{2}}{2} & -\frac{\sqrt{2}}{2} 
\end{bmatrix}
\quad \text{e} \quad 
R = 
\begin{bmatrix} 
\sqrt{2} & \frac{\sqrt{2}}{2} \\[6pt] 
0 & \frac{\sqrt{2}}{2} 
\end{bmatrix}
\]
\end{frame}

% Slide 14
\begin{frame}{Aplicações}
\begin{itemize}
    \item Resolução de sistemas lineares.
    \item Computação gráfica.
    \item Processamento de sinais.
    \item Inteligência Artificial.
\end{itemize}
\end{frame}

% Slide 15
\begin{frame}{Vantagens}
\begin{itemize}
    \item Estabilidade numérica.
    \item Eficiência computacional.
    \item Facilidade algébrica.
    \item Aplicação em grandes sistemas.
\end{itemize}
\end{frame}

% Slide 16
\begin{frame}{Desvantagens}
\begin{itemize}
    \item Alto custo computacional em matrizes grandes.
    \item Gram-Schmidt clássico pode gerar erros numéricos.
    \item Necessidade de precisão computacional.
\end{itemize}
\end{frame}

% Slide 17
\begin{frame}{QR na Computação}
Bibliotecas utilizadas:
\begin{itemize}
    \item MATLAB
    \item NumPy
    \item SciPy
    \item Octave
\end{itemize}
\end{frame}

% Slide 18
\begin{frame}{Conclusão}
\begin{itemize}
    \item QR é essencial em Álgebra Linear.
    \item Possui ampla aplicação prática.
    \item Gram-Schmidt é base da decomposição.
    \item Muito importante para computação científica.
\end{itemize}
\end{frame}

% Slide 19
\begin{frame}{Obrigado!}
\centering

\Huge Dúvidas?

\vspace{1cm}

\Large Obrigado pela atenção!
\end{frame}

\end{document}
